\section*{Chapitre \ref{ch:b3:exploration-seismology} --- Sismique d'exploration}

\begin{solution}{exo:b3:exploration-seismology:1}
$Z = \rho v$~: craie $2000\times2500 = \qty{5.00e6}{kg.m^{-2}.s^{-1}}$, argile
$2300\times2300 = \qty{5.29e6}{kg.m^{-2}.s^{-1}}$, calcaire $2600\times4500 =
\qty{11.70e6}{kg.m^{-2}.s^{-1}}$. Craie--argile~: $R = (5.29 - 5.00)/(5.29 + 5.00)
= 0.29/10.29 = 0.0282$~; argile--calcaire~: $R = (11.70 - 5.29)/(11.70 + 5.29) =
6.41/16.99 = 0.377$. L'énergie réfléchie est $R^2 = 0.142$~: le sommet du
calcaire renvoie \qty{14.2}{\percent} de l'énergie incidente, l'interface craie--argile
en renvoie moins de \qty{0.1}{\percent}.
\end{solution}

\begin{solution}{exo:b3:exploration-seismology:2}
$t_0 = 2\times800/2000 = \qty{0.800}{s}$~; $t = \sqrt{0.800^2 + (1000/2000)^2} =
\sqrt{0.89} = \qty{0.9434}{s}$~; le retard vaut \qty{0.1434}{s}. Par
l'approximation donne $1000^2/(2\times2000^2\times0.800) = \qty{0.1563}{s}$, trop
grand de \qty{9}{\percent}~: pour un déport supérieur à la profondeur, la relation $x \ll vt_0$ ne
est plus vérifiée.
\end{solution}

\begin{solution}{exo:b3:exploration-seismology:3}
Les points milieux se déplacent de la moitié de l'intervalle entre récepteurs~: l'intervalle entre points milieux communs est
$12.5/2 = \qty{6.25}{m}$. L'\omterm{def:b3:exploration-seismology:cmp}{ordre de couverture} vaut $240\times12.5/(2\times25) = 60$. Un \omterm{def:b3:exploration-seismology:cmp}{ordre de
couverture} de 120 exige un tir tous les $240\times12.5/(2\times120) = \qty{12.5}{m}$.
\end{solution}

\begin{solution}{exo:b3:exploration-seismology:4}
$\lambda = v/f$~: \qty{100}{m} dans le grès, $\lambda/4 = \qty{25}{m}$~;
\qty{150}{m} dans le calcaire, $\lambda/4 = \qty{37.5}{m}$. Avec
$\lambda = 2800/30 = \qty{93.3}{m}$ et $z = \qty{1000}{m}$,
$r_F \approx \sqrt{\lambda z/2} = \sqrt{46\,667} = \qty{216}{m}$ (\qty{217}{m}
en ajoutant le terme $\lambda^2/16$).
\end{solution}

\begin{solution}{exo:b3:exploration-seismology:5}
\omterm{thm:b3:exploration-seismology:dix}{Équation de Dix}~: $v^2 = (2800^2\times0.700 - 2400^2\times0.500)/0.200 = (5\,488\,000 -
2\,880\,000)/0.200 = \qty{1.304e7}{m^2.s^{-2}}$, donc $v = \qty{3611}{m.s^{-1}}$.
La couche est traversée en \qty{0.100}{s} dans un sens~: $h = 3611\times0.100 =
\qty{361}{m}$.
\end{solution}

\begin{solution}{exo:b3:exploration-seismology:6}
(a) Soit $d$ la distance du point milieu M au réflecteur, de sorte que
$t_0 = 2d/v$. Le rai réfléchi de la source S au récepteur G est aussi long
que la droite reliant G à l'image S$'$ de S par rapport au réflecteur. Plaçons M à
l'origine, la droite suivant le pendage~; S$'$ se trouve à la distance $2d_S$ de S le long
de la normale, avec $d_S = d - (x/2)\sin\varphi$. Le carré de S$'$G vaut
$x^2 + 4d_S^2 + 4d_Sx\sin\varphi = x^2 + 4d^2 - x^2\sin^2\varphi$, tous les termes
en $dx$ s'annulant, donc $v^2t^2 = 4d^2 + x^2\cos^2\varphi$ et
$t^2 = t_0^2 + x^2\cos^2\varphi/v^2$. (b) La collection de traces suit une hyperbole de vitesse
$V = v/\cos\varphi = 2000/\cos\ang{12} = \qty{2045}{m.s^{-1}}$~: un pendage rend la
\omterm{def:b3:exploration-seismology:velocities}{vitesse de sommation} trop élevée, ici de \qty{2.2}{\percent}.
\end{solution}

\begin{solution}{exo:b3:exploration-seismology:7}
Par $\sqrt{60} = 7.75$. Doubler le rapport signal sur bruit exige quatre fois
autant de traces~: un \omterm{def:b3:exploration-seismology:cmp}{ordre de couverture} de 240.
\end{solution}

\begin{solution}{exo:b3:exploration-seismology:8}
La profondeur tracée $vt/2$ varie de $\qty{1.5}{km.s^{-1}}\times\qty{0.40}{s}
= \qty{0.60}{km}$ par kilomètre~: $\tan\theta_a = 0.60$, $\theta_a = \ang{31.0}$.
Par le théorème du déport nul $\sin\theta = 0.60$~: $\theta = \ang{36.9}$.
\end{solution}

\begin{solution}{exo:b3:exploration-seismology:9}
{\sloppy Chaque pic place une copie de $w$ centrée sur lui~: $0.3\,w$ sur les échantillons 0--2 donne
$(-0.15, 0.3, -0.15)$, et $-0.2\,w$ sur les échantillons 2--4 donne $(0.1, -0.2, 0.1)$.
Leur somme vaut $(0, -0.15, 0.3, -0.05, -0.2, 0.1, 0)$, écrite de l'échantillon
précédant la série jusqu'à celui qui la suit. Les réflexions culminent aux échantillons 1 et 3, avec
leurs vrais signes~; elles se superposent à l'échantillon 2, où $-0.15 + 0.1 = -0.05$~: une
légère interférence, mais deux pics distincts.\par}
\end{solution}

\begin{solution}{exo:b3:exploration-seismology:10}
$\arcsin(1800/2900) = \ang{38.37}$ et $\arcsin(1800/4400) = \ang{24.15}$, donc
$i_c = (38.37 + 24.15)/2 = \ang{31.26}$ et $\theta = (38.37 - 24.15)/2 =
\ang{7.11}$~; $v_2 = 1800/\sin\ang{31.26} = \qty{3469}{m.s^{-1}}$. La moyenne
harmonique $2\times2900\times4400/(2900 + 4400) = \qty{3496}{m.s^{-1}}$ est
supérieure de \qty{0.8}{\percent}~; multipliée par $\cos\theta = 0.9923$, elle donne
\qty{3469}{m.s^{-1}}, l'approximation usuelle pour les faibles pendages
$v_2 \approx \cos\theta\,2v_dv_u/(v_d + v_u)$.
\end{solution}

\begin{solution}{exo:b3:exploration-seismology:11}
(a) Le signal composite vaut $R[w(t) - w(t - \Delta t)] \approx R\,\Delta t\,w'(t)$ par
le développement de Taylor au premier ordre. (b) Son amplitude est proportionnelle à
$\Delta t$, donc un banc plus mince que l'\omterm{def:b3:exploration-seismology:resolution}{épaisseur d'accord} renvoie encore un signal,
de forme $w'$, mais son toit et sa base ne peuvent être pointés comme deux événements~: il
est détecté, mais non résolu. Si $R$ et l'ondelette sont connus, par exemple grâce à un
forage, l'amplitude donne $\Delta t$ et donc l'épaisseur. (c) L'amplitude est
proportionnelle à l'épaisseur~: le rapport vaut $5/10 = 1/2$.
\end{solution}

\begin{solution}{exo:b3:exploration-seismology:12}
(a) $t(x) = (2/v)\sqrt{z^2 + x^2}$~: $t(0) = 2\times600/3000 = \qty{0.400}{s}$ et
$t(300) = (2/3000)\sqrt{600^2 + 300^2} = \qty{0.4472}{s}$. (b) $t^2 = t_0^2 +
x^2/(v/2)^2$, la forme NMO avec la vitesse $v/2$. Une réflexion de même temps au sommet
sur une section à déport nul, ne montrerait aucun retard~; comparée à une
hyperbole NMO de même $t_0$ dans une collection, le retard de la
$x^2/(2(v/2)^2t_0)$ est quatre fois plus grand, car le temps dépend deux fois de la
distance au diffuseur deux fois, à l'aller et au retour. (c) La migration somme la
section le long de cette hyperbole et place l'énergie à son sommet~: la diffraction
se réduit à un point en $x = 0$, $z = \qty{600}{m}$.
\end{solution}

\begin{solution}{pb:b3:exploration-seismology:1}
\textbf{1.} Depuis O, $400/(486.1 - 425.2) = 400/0.0609 = \qty{6568}{m.s^{-1}}$~;
depuis E, $400/(486.1 - 372.4) = 400/0.1137 = \qty{3518}{m.s^{-1}}$. L'onde
conique issue de O semble rapide~: O est l'extrémité aval-pendage, et ses tirs
remontent le pendage~; le réfracteur s'élève de O vers E.

\textbf{2.} $\arcsin(2400/3518) = \ang{43.02}$, $\arcsin(2400/6568) =
\ang{21.43}$~: $i_c = \ang{32.22}$, $\theta = \ang{10.79}$.

\textbf{3.} $v_2 = 2400/\sin\ang{32.22} = \qty{4501}{m.s^{-1}}$~: dans la
plage du calcaire, le sommet du calcaire jurassique.

\textbf{4.} Les temps d'interception valent $t - x/v_{\text{app}}$ à \qty{1200}{m}~:
$t_{i,W} = 0.4861 - 1200/6568 = \qty{0.3034}{s}$ et $t_{i,E} = 0.4861 -
1200/3518 = \qty{0.1450}{s}$. Avec $\cos i_c = 0.8460$, $h = v_1t_i/(2\cos
i_c)$~: $h_W = 2400\times0.3034/1.692 = \qty{430.4}{m}$ et $h_E =
\qty{205.7}{m}$.

\textbf{5.} En divisant par $\cos\theta = 0.9823$~: \qty{438.1}{m} sous O et
\qty{209.4}{m} sous E. Alors $(438.1 - 209.4)/1200 = 0.1906 =
\tan\ang{10.79}$, le pendage de la question~2.

\textbf{6.} Le schiste ne renvoie aucune onde conique~: il est invisible, et la
campagne le compte avec la craie. Ici $v_1 = \qty{2.40}{km.s^{-1}}$ était donnée
comme la moyenne des deux~; si la vitesse de l'onde directe dans la craie, \qty{2.50}{km.s^{-1}},
avait été utilisée pour toute la couverture, chaque profondeur serait sortie trop grande.

\textbf{7.} $V = 600/\sqrt{0.5437^2 - 0.4878^2} = 600/0.2401 =
\qty{2499}{m.s^{-1}}$ et $V = 600/\sqrt{0.6240^2 - 0.5748^2} =
\qty{2471}{m.s^{-1}}$. La seconde moyenne la craie avec le schiste, plus lent,
situé dessous~: la moyenne baisse.

\textbf{8.} \qty{2811}{m.s^{-1}} et \qty{2829}{m.s^{-1}}.

\textbf{9.} Couche 1~: \qty{2499}{m.s^{-1}}. Couche 2~: $v^2 = (2471^2\times0.5748
- 2499^2\times0.4878)/0.0870$, $v = \qty{2308}{m.s^{-1}}$. Couche 3~:
$v^2 = (2811^2\times0.6548 - 2471^2\times0.5748)/0.0800$,
$v = \qty{4561}{m.s^{-1}}$. Couche 4~: $v^2 = (2829^2\times0.7348 -
2811^2\times0.6548)/0.0800$, $v = \qty{2972}{m.s^{-1}}$.

\textbf{10.} Craie (\qty{2499}{m.s^{-1}}), schiste (\qty{2308}{m.s^{-1}}),
calcaire (\qty{4561}{m.s^{-1}}), grès (\qty{2972}{m.s^{-1}}), chacun dans sa
plage publiée. Les réflexions sont l'interface craie--schiste, le sommet
du calcaire, l'interface calcaire--grès et la discordance au
sommet du socle schisteux.

\textbf{11.} La craie et le schiste ont presque la même impédance~: $R = 0.028$
(exercice~1), contre 0.38 pour le sommet du calcaire.

\textbf{12.} L'hyperbole de vitesse $V_{\text{rms}}$ ne tient qu'au second ordre
en déport~; pour des déports de \qty{2}{km}, soit le double de la profondeur du
réflecteur le plus profond, les temps s'en écartent, et une hyperbole ajustée
sur tout le dispositif donnerait de fausses vitesses, et des intervalles de Dix encore pires. Les grands déports
étirent aussi les ondelettes après correction.

\textbf{13.} Temps doubles dans les couches~: \qty{0.4878}{s}, \qty{0.0870}{s},
\qty{0.0800}{s}, \qty{0.0800}{s}. Épaisseurs $v\,\Delta t/2$~: \qty{609.5}{m},
\qty{100.4}{m}, \qty{182.4}{m}, \qty{118.9}{m}.

\textbf{14.} Profondeurs sous le sol~: \qty{609.5}{m}, \qty{709.9}{m},
\qty{892.3}{m}, \qty{1011.2}{m}. Le schiste, le calcaire et le grès sortent
à moins de \qty{0.4}{m}, \qty{2.4}{m} et \qty{1.1}{m} de leurs épaisseurs, l'erreur
croissant dans le calcaire rapide et mince, où une petite erreur de temps coûte
le plus.

\textbf{15.} $\bar v = 2\times1011.2/0.7348 = \qty{2752}{m.s^{-1}}$. La \omterm{def:b3:exploration-seismology:velocities}{vitesse
quadratique moyenne} pondère chaque couche par $v^2$ et n'est jamais inférieure à
la moyenne pondérée par le temps~: $2829\times0.7348/2 = \qty{1039}{m}$, soit \qty{28}{m} de trop en profondeur.

\textbf{16.} $\lambda = 2308/30 = \qty{76.9}{m}$, $\lambda/4 = \qty{19.2}{m}$~;
l'estimation de Ricker donne $2308/(5.2\times30) = \qty{14.8}{m}$. Le schiste
fait \qty{100}{m} d'épaisseur, soit cinq fois son \omterm{def:b3:exploration-seismology:resolution}{épaisseur d'accord}~: son sommet et sa base
sont résolus. Il s'accorderait, selon la règle du quart de longueur d'onde, à $f = v/(4b) =
2308/400 = \qty{5.77}{Hz}$.

\textbf{17.} $\lambda = 2829/30 = \qty{94.3}{m}$, $z = \qty{1011}{m}$~:
$r_F \approx \sqrt{94.3\times1011/2} = \qty{218}{m}$ (\qty{220}{m} avec le
terme $\lambda^2/16$).

\textbf{18.} $t(x)^2 = t_0^2 + (2x/v)^2$, donc $v = 2x/\sqrt{t^2 - t_0^2} =
400/\sqrt{0.5732^2 - 0.550^2} = 400/0.1614 = \qty{2478}{m.s^{-1}}$.

\textbf{19.} $z = vt_0/2 = 2478\times0.550/2 = \qty{681}{m}$.

\textbf{20.} La profondeur tracée $vt/2$ varie de $1.24\times0.167 = 0.207$ par
unité de longueur~: $\theta_a = \arctan0.207 = \ang{11.70}$. Le vrai pendage
vérifie $\sin\theta = 0.207$~: $\theta = \ang{11.95}$, le pendage du flanc.

\textbf{21.} $d = 2480\times0.655/2 = \qty{812}{m}$~; la migration déplace la
réflexion de $d\sin\theta = 812\times0.207 = \qty{168}{m}$ vers le haut du
pendage, en direction de la crête de l'anticlinal, et la fait remonter de \qty{812}{m} à $d\cos\theta =
\qty{795}{m}$.

\textbf{22.} $2500\times0.160/2 = \qty{200}{m}$.

\textbf{23.} D'après les questions 13 et 14, la discordance de socle se trouve
\textbf{\qty[text-series-to-math=true]{1011}{m}} sous le sol au \omterm{def:b3:exploration-seismology:cmp}{point milieu commun} S.
\end{solution}
